{\begingroup
\renewcommand{\hom}{\operatorname{Hom}}
\newcommand{\tM}{\widetilde{M}}
\newcommand{\ann}{\operatorname{Ann}}
\newcommand{\ord}{\operatorname{ord}}
\newcommand{\supp}{\operatorname{supp}}
\newcommand{\fun}{\operatorname{Funct}}
\newcommand{\sch}{\operatorname{Sch}}
\newcommand{\spec}{\operatorname{Spec}}
\newcommand{\proj}{\operatorname{Proj}}
\newcommand{\ring}{\operatorname{Ring}}
\newcommand{\grp}{\operatorname{Ab}}
\newcommand{\pic}{\operatorname{Pic}}
\newcommand{\qcoh}{\operatorname{Qcoh}}
\newcommand{\op}{\text{op}}
\newcommand{\pt}{\text{pt}}
\newcommand{\g}{\mathscr{G}}
\renewcommand{\l}{\mathscr{L}}
\newcommand{\f}{\mathscr{F}}
\renewcommand{\o}{\mathcal{O}}
\renewcommand{\P}{\mathbb{P}}
\newcommand{\p}{\mathfrak{p}}
\newcommand{\m}{\mathfrak{m}}
\newcommand{\C}{\mathbb{C}}
\newcommand{\Z}{\mathbb{Z}}
\newcounter{paranum} 
\newcommand{\np}{\stepcounter{paranum}\noindent\textbf{\arabic{paranum}.}\quad}

\chapter{Algebraic Geometry}
Let us review concepts from  Algebraic Geometry.

\np Let $A\in \ring$, then we have $\sch_{/{\spec(A)}}\simeq \ring_{A/}$.
The correspondence is given by $(X\mapsto \spec(A))\leadsto A\to \Gamma(X,\o_X)$.

\np If we consider $X\mapsto \mathbb{A}^1_{A}$, it corresponds to $A[t]\mapsto \Gamma(X,\o_X)$ which uniquely determined by the image of $t$.
In this sense, we obtain that $(X\mapsto \mathbb{A}^1_{A})\leadsto \Gamma(X,\o_X)$. 
Moreover, we have $\hom_{S}(X,\mathbb{A}_S^1)\simeq \hom_{\spec(\mathbb{Z})}(X,\mathbb{A}^1)\simeq \Gamma(X,\o_X)$.

\np   Let $A$ be a ring, we have $\proj_{/\spec(A)}\simeq \text{Mod}_A$.


\np Structure sheaf $O_X:X^\op\to \ring$.

\np Let $j_!$ denotes extension by zero and $U$ is an open subscheme of $X$ then we obtain that $\Gamma(U, \mathcal{F}) \cong \text{Hom}_{\mathcal{O}_X}(j_! \mathcal{O}_U, \mathcal{F})$.

\np Serre’s criterion for Normality: A Noetherian ring $A$ is normal
if and only if it has properties R1 and S2.
R1 means $A_\mathfrak{p}$ is regular(smooth) for any codimension one prime ideal $\p$. 
S2 means $\text{depth}(A_\p)=\min(2,\dim(A_\p))$ for any prime ideal.

\np $\text{depth}(A_\p)$ means the maximal length of regular sequence $(a_1,a_2,\dots,)$ in $\p$.
Regular means that $a_i$ is not a zero divsor in $A/(a_1,\dots,a_i)$.

\np Sheaf of $O_X$-module:
\[
\begin{tikzcd}[column sep=large, row sep=large]
& \mathcal{M}od \arrow[d, "P"] \\
\text{Open}(X)^{\text{op}} \arrow[ur, "\mathcal{F}"] \arrow[r, "\mathcal{O}_X"'] & \text{Ring}
\end{tikzcd}
\]

\np $f:X\to Y$, then $f^*\g  = f^{-1}\g \otimes_{f^{-1}\o_Y}\o_X $

\np Let $M$ be $A$-module and $X=\spec(A)$. Then $\widetilde{M}:=\underline{M}\otimes \o_X$ where $\underline{M}$ denotes the constant sheaf of $M$.

\np Ideal sheaf is the kernal of $i^\#:\o_X\to i_*\o_Y$ induced by closed immersion $i:Y\to X$.

\np We call a sheaf $\f$ is quasi-coherent if for any affine open subscheme $U$, there exist an $\o_X(U)$-module $M$ such that $\f|_{U}\simeq \widetilde{M}$.
Hence, we call it coherent if $X$ is Noetherian.

\np Let $X$ be any affine scheme and $\f$ be an arbitrary quasi-coherent sheaf. Then we obtain that $H^1(X,\f)=0$.
It means that if we consider the exact sequence of ideal sheaf 
\begin{align*}
    0 \to \mathcal{I}_Y \to \mathcal{O}_X \to i_* \mathcal{O}_Y \to 0
\end{align*}
The global section functor is exact 
\begin{align*}
    0 \to \Gamma(X, \mathcal{I}_Y) \to \Gamma(X, \mathcal{O}_X) \to \Gamma(X, i_* \mathcal{O}_Y)\to H^1(X, \mathcal{I}_Y) = 0.
\end{align*}
Equivalently, we say any regular function on closed subscheme of $X$ can be extended to $X$.

\np For any scheme $X$, S2 condition is equivalent to $j_*(\o_U)\simeq \o_X$ for any $\text{codim}(X\setminus U)\ge 2$

\np For open subscheme $U$ of $X$, there is an exact sequence
\begin{align*}
    0 \to i_! \mathcal{O}_U \to \mathcal{O}_X \to j_* \mathcal{O}_Z \to 0
\end{align*}
where $i_!\o_U$ is the push-forward with compact support. The inclusion is due to the extension by zero.

\np We also have exact sequence 
$$0 \to \mathcal{H}^0_Z(\mathcal{O}_X) \to \mathcal{O}_X \to i_* \mathcal{O}_U \to \mathcal{H}^1_Z(\mathcal{O}_X) \to \dots$$
where $\o_X\to i_*\o_U$ is restriction. $\mathcal{H}^0_Z(\mathcal{O}_X)$ means the sheaf of supporting on $Z$. 
It also can be defined by the sheafication of presheaf functor $U\mapsto \Gamma_{Z}(X,\o_X)\cap U$ where $\Gamma_Z$ means supporting on $Z$.

\np In the projective case, we have an similar exact sequence. Given that an projecitve shceme $X=\proj(S)$, then we interpret it as the open subscheme of $CX$ excluding zero.
The exact sequence follows 
\begin{align*}
    0\to H^0_{\mathfrak{m}}(S) \to \Gamma(CX,\o_{CX})= S \to \Gamma_*(\o_X)\to  H^1_{\mathfrak{m}}(S) \to H^1(CX,\o_{CX})=0.
\end{align*}

\np Let $M$ be $R$-module, we define $\text{Ass}(M)=\{\text{Ann}_R(m)\}_{m\in M}\cap \spec(R)$. Then $\text{Supp}(M)=\bigcup_{\p\in \text{Ass}(M)}V(\p)$.
To see that we consider if $\p=\text{Ann}_R(m)$ then $M_\p\neq 0$.

\np Let $X=\spec(R)$ and $f\in R$, then $\text{Supp}(f)=V(\text{Ann}_R(f))$. Hence, if $f$ is not zero divsor, then $\text{Supp}(f)=X$.

\np If we consider $_R R$, then $\text{Ass}(R)=\{\sqrt{\mathfrak{q}_1},\dots,\sqrt{\mathfrak{q}_n}\}$
where $(0) = {\mathfrak{q}_1}\cap \dots\cap {\mathfrak{q}_n}$ be the primary decomposition of zero ideal.

\np Suppse that $I$ is radical and $Z=V(I)$. We have $\mathcal{H}^0_Z(\mathcal{O}_X)=\{x\in\o_X\mid \sqrt{\text{Ann}_R(x)}\supseteq I\}$.
Hence, we can obtain that $\mathcal{H}^0_Z(\mathcal{O}_X)=\bigcap_{\p_i\not\supseteq I}\mathfrak{q}_i$. 


\np For a graded ring $S$ with $M$ is a graded $S$-module. Define $\widetilde{M}(D_+(f)):=\{m/f^n|\deg(m)=n\}$.

\np Let $S$ be a graded ring, $S(n)$ denotes twisted module defined by $S(n)_{d}=S_{n+d}$. Define the twisted sheaf $\o_X(n):=\widetilde{S(n)}$ and twisted sheaf of module 
$\f(n):=\f\otimes_{\o_X}\o_X(n)$.


\np Let $S$ be the graded ring generated by $S_1$ as $S_0$-algebra. Let $X=\proj(S)$.
Then
\begin{itemize}
    \item $\o_X(n)$ is an invertible sheaf on $X$;
    \item $\widetilde{M}(n)\simeq \widetilde{M(n)}$;  
    \item $\o_X(n)\otimes\o_X(m)\simeq \o_X(n+m)$
\end{itemize}

\np The section functor is given by $\Gamma_*(\f)= \bigoplus_{n\in \mathbb{Z}}\Gamma(X,\f(n))$. It is right adjoint functor of $\widetilde{(\cdot)}$.
Hence if $S$ is a graded ring generated by $S_0$-algebra $S_1$ and $\f$ is quasi-coherent, then $\Gamma_*(\f)\simeq \f$.

\np Extension of Sections: Let $X$ be a scheme and $\l\in \pic(X),\f\in\qcoh(X)$. 
Suppose $X$ is quasi-compact and has a finite open cover $\{U_i\}$ such that $\l|_{U_i}$ is free and $U_i\cap U_j$ is quasi compact.
Given that $t\in\Gamma(X_f,\f)$, then for some $n>0$, the section $f^nt\in\Gamma(X_f,\f\otimes \l^{\otimes n})$ can be extended to a global section of $\f\otimes \l^{\otimes n}$.

\np we call a morphism $f:X\to Y$ is projective if $f$ can be factorized as form $X\to \mathbb{P}_{Y}^n\to Y$
where $X\to\mathbb{P}_Y^n$ is a closed immersion. 
In this case, we call $X$ is a projective scheme over $Y$.  If $Y=\spec(A)$, then exists a graded ring $S$ finitely generated by $A=S_0$-algebra $S_1$.

\np we call a morphism is quasi projective if $X\to\mathbb{P}_Y^n$ is a locally closed immersion which means there exists closed subscheme $Z$ such that $X\to Z$ is an open immersion.  

\np For any scheme $Y$, define twisting sheaf $\o(d)$ on $\mathbb{P}_Y^r$ by $g^*(\o(d))$ where $g:\mathbb{P}_Y^r \to \mathbb{P}_{\mathbb{Z}}^r$.

\np $X$ is a shceme over $Y$. We call invertible sheaf $\l$ is very ample relative to $Y$ if there is a locally closed immersion $i:X \to \mathbb{P}_Y^r$ for some $r$ such that $i^*(\o(1))\simeq \l$. 


\np We call $\f$ is generated by global sections if there is a family of global sections $\{s_i\}$ such that $\f_x$ is generated by $\{s_i(x)\}$.

\np Let $A$ be a finitely generated $k$-algebra with graded structure. Let $X$ be projecitve over $\spec(A)$ and $\f$ is an arbitrary coherent sheaf. Then $\Gamma(X,\f)$ is a finitely generated $A$-module.

\np Let $Y$ is quasi-projecitve over $\spec(A)$. Equivalently, $Y$ is a quasi-projecitve subscheme of $\mathbb{P}_{A}^n$. Then we have a very ample line bundle o



\np Flasque sheaf is a sheaf such that for any two $V\subseteq U$, we have 
\begin{align*}
    \rho_{U,V}:U\to V
\end{align*}
is surjective. An example is consider $i:\pt \to X$ and the sheaf $i_*(\o_\pt)$ is a Flasque sheaf of $X$.

\np Every finite type affine $k$-scheme is quasi-projective.

\np Let $X=\spec R$ with a sheaf of module $\tM$. Then we have 
\begin{align*}
    V(\ann_R(M)) = \supp_X(\tM).
\end{align*}
Let $R=k[x]$ and $R=k[x]/(x^2)$. Then the support of $M$ is given by 
\begin{align*}
    Y = \spec(k[x]/(x^2))
\end{align*}
But its topological support is $Y_{\text{red}}$.

\np Let $B = A[x_1,\cdots,x_n]/I$. The Kähler differential is given by 
\begin{align*}
    \Omega_{B/A} = \left(\bigoplus_{i=1}^n B dx\right)/\langle df,f\in I \rangle.
\end{align*}



\np Let $A\to B\to C$ be rings. There First exact sequence of Kähler differential is 
\begin{align*}
    \Omega_{B/A}\otimes_B C \to \Omega_{C/A} \to  \Omega_{C/B} \to 0
\end{align*}

Actually, we have 
\begin{align*}
    \Omega_{B/A} = \Omega_{B/\C}/\langle dg,g\in A \rangle 
\end{align*}

If we let $C=B/I$ then we obtain a second exact sequence 
\begin{align*}
    I/I^2 \xrightarrow{\delta} \Omega_{B/A}\otimes_B C \to \Omega_{C/A} \to 0.
\end{align*}

Where $I/I^2$ is conormal module. The normal module is the dual of conormal module which define by 
\begin{align*}
    N_{C/A} = \hom_{C}(I/I^2,C).
\end{align*}

The tangent bundle is given by 
\begin{align*}
    T_{C/A} = \hom_{C}(\Omega_{C/A},C).
\end{align*}

Let $X=\spec(A)$ and $Y=\spec(A/I)$ and smooth.
We have the exact sequence 
\begin{align*}
    0\to N^\vee _Y \xrightarrow{\delta} \Omega_{X}\otimes_{\o_X} \o_Y \to \Omega_{Y} \to 0.
\end{align*}
where $N^\vee _{Y/X}$ denotes the sheaf induced by conormal module $I/I^2$. Dualize we then have 
\begin{align*}
    0\to T_Y \to  T_X\otimes_{\o_X} \o_Y \to N_{Y/X} \to 0.
\end{align*}

\[
X=\mathbb A^2_{\mathbb C}=\spec\mathbb C[x,y],\quad Y=V(y-x^2),\quad I=(y-x^2),\quad \mathcal O_Y\simeq\mathbb C[x].
\]
\[
\Omega_X|_Y=\mathcal O_Ydx\oplus\mathcal O_Ydy,\qquad 
\Omega_Y\simeq(\mathcal O_Ydx\oplus\mathcal O_Ydy)/(dy-2x\,dx),\qquad T_Y=\mathcal O_Y\left(\partial_x+2x\partial_y\right).
\]
\[
I/I^2\simeq d(y-x^2)=dy-2x\,dx. \quad 
N_{Y/X}\simeq
(\mathcal O_Y\partial_x\oplus\mathcal O_Y\partial_y)/
\mathcal O_Y(\partial_x+2x\partial_y).
\]

\np In local we also have an exact sequence
\begin{align*}
    0\to \m_x/\m_x^2 \to \o_{X,x} \to \o_{X,x}/{\m_x} \to 0 
\end{align*}
where $\m_x/\m_x^2$ is fiber of conormal sheaf.

\np We define canonical sheaf of $X$ by $\omega_X=\bigwedge^n\Omega_{X/k}$.
If $X$ is normal(integrally closed) let $\eta\in \omega_X$, we define canonical divsor by 
\begin{align*}
    K_X = \sum \ord_D(\eta) \cdot D.
\end{align*}
Then $\o_X(K_X) = w_X$.

\np Let $H=V(x_0)$ be a hyperplane of $\P^n$. Then $\pic(\P^n)=\langle H \rangle$.

\np Consider $\o_{\P^{2}}(-1)$ as tautology line bundle. 

We actually have 
\begin{align*}
    \o(-1)(U_0) = \C[\frac{x_1}{x_0},\frac{x_2}{x_0}]\cdot x_0^{-1} = \C[x_1,x_2]
\end{align*}

\np Let $Y=\spec A$ and let $X=\P^n_A$. Then there is an exact sequence of sheaves.
\begin{align*}
    0\to \Omega_{X/Y} \to \o_X(-1)^{\oplus (n+1)}\to \o_X\to 0.
\end{align*}

We can check by take $Y=\spec\C$ and $X=\P^2_\C$. Consider the $U_0$, then we obtain that 
\begin{align*}
    \Omega_{X/Y}(U_0) \simeq \C[\frac{x_1}{x_0},\frac{x_2}{x_0}] d(x_1/x_0)\oplus \C[\frac{x_1}{x_0},\frac{x_2}{x_0}] d(x_2/x_0) \\
    \o_X(-1)^{\oplus 3}(U_0) \simeq \bigoplus_{i=1}^3 \C[\frac{x_1}{x_0},\frac{x_2}{x_0}]\cdot x_0^{-1}\\
\end{align*}



\endgroup}

